% Generated draft: requires mathematical and source review.
\documentclass[11pt]{article}
\usepackage[margin=25mm]{geometry}
\usepackage{fontspec}
\setmainfont{Latin Modern Roman}
\usepackage{amsmath,amssymb,mathtools,unicode-math}
\setmathfont{Latin Modern Math}
\usepackage{enumitem,needspace,xurl,hyperref}
\setlength{\parindent}{0pt}
\setlength{\parskip}{5pt}
\setlength{\emergencystretch}{3em}
\newcommand{\N}{\mathbb{N}}
\newcommand{\Z}{\mathbb{Z}}
\newcommand{\Q}{\mathbb{Q}}
\newcommand{\R}{\mathbb{R}}
\newcommand{\C}{\mathbb{C}}
\newcommand{\F}{\mathbb{F}}
\newcommand{\A}{\mathbb{A}}
\newcommand{\PP}{\mathbb{P}}
\newcommand{\E}{\mathbb{E}}
\newcommand{\eps}{\varepsilon}
\newcommand{\dd}{\,\mathrm{d}}
\begin{document}
\section*{UnsolvedMath ID: 30001804}
\textbf{Source problem code:} catalogue entry
\subsection{Definitions and mathematical statement}
Write \(G_g=\operatorname{Mod}(\Sigma_g^1)\) and \(M_g=\operatorname{St}(\Sigma_g^1)\) for the mapping class group and surface Steinberg module named in the source. A mapping class group consists of isotopy classes of the allowed self-homeomorphisms of the surface, with multiplication induced by composition. The supplied record does not specify whether the superscript \(1\) denotes a marked point, a puncture, or a boundary component, nor the corresponding requirements on homeomorphisms and isotopies. It also does not specify the intended genus range or the coefficient ring and homological definition of the Steinberg module. These conventions must be recovered from the cited sources before a fully specified surface-theoretic statement can be given.

The requested presentation can nevertheless be expressed precisely in algebraic terms, conditional on fixing those conventions. Let \(R\) be the coefficient ring used in the definition of \(M_g\), and regard the mapping class group action as a left action, so that \(M_g\) is a left module over the group ring \(R[G_g]\). Elements of \(R[G_g]\) are finite formal sums of group elements with coefficients in \(R\), and their action on \(M_g\) extends the group action linearly. A single module generator is an element \(s_g\in M_g\) satisfying \(M_g=R[G_g]s_g\). The supplied record reports that such a generator is known.

\textbf{Requested statement.} For each genus in the intended range, describe the \(G_g\)-module structure of \(M_g\) by giving a presentation using its single module generator. More explicitly, identify a cyclic generator \(s_g\), or specify the known generator being used, and give an explicit family of relations \(\{r_{g,\lambda}\}_{\lambda\in\Lambda_g}\subseteq R[G_g]\) such that the map
\[
\Phi_g:R[G_g]\longrightarrow M_g,\qquad a\longmapsto a\cdot s_g,
\]
is surjective and
\[
\ker\Phi_g
=\operatorname{Ann}_{R[G_g]}(s_g)
=\sum_{\lambda\in\Lambda_g}R[G_g]r_{g,\lambda}.
\]
Equivalently, give an explicit isomorphism of left \(R[G_g]\)-modules
\[
M_g\cong R[G_g]\Big/\sum_{\lambda\in\Lambda_g}R[G_g]r_{g,\lambda},
\qquad [1]\longmapsto s_g.
\]
Here the sum consists of finite sums of left multiples of the displayed relations. Merely naming the annihilator does not provide the requested explicit presentation. The record does not explicitly require \(\Lambda_g\) to be finite; whether a finite or uniformly described family is intended remains unresolved. Taking \(R=\Z\) would give the integral formulation, but that coefficient convention is not established by the supplied record.
\subsection{Short English statement}
Describe the surface Steinberg module under the mapping class group action by explicit relations on one module generator. The surface convention, genus range, coefficients, and precise Steinberg module definition require confirmation from the cited sources. The supplied dated assessment labels the general problem open and reports that Broaddus obtained a finite presentation with several generators and a cyclic generator, with explicit conversion to a presentation on one generator only in low genus.
\subsection{Sources}
\begin{itemize}
\item Local UnsolvedMath registry: supplied complete record for permanent ID \(30001804\), source problem code \texttt{catalogue entry}, entitled \emph{Mapping-Class-Group Presentations of Surface Steinberg Modules}. The record gives the proposed year as 2011 and names no individual proposer.
\item \emph{Arithmetic Groups vs. Mapping Class Groups: Similarities, Analogies and Differences} (2011), Oberwolfach Reports, pp.\ 1637--1708, DOI \texttt{10.4171/owr/2011/30}. Supplied source report and literature-source link labelled ``original Oberwolfach report'': \url{https://doi.org/10.4171/owr/2011/30}. The record describes these reports as documenting mathematical workshops and their open-problem discussions.
\item Supplied literature-source link labelled ``Broaddus--Steinberg module'': \url{https://arxiv.org/abs/0711.0011}. The supplied assessment attributes a finite presentation with several generators and a cyclic generator to Broaddus; it says that the requested general presentation on the single generator is not supplied by that work and that only low-genus cases are explicitly converted.
\item Literature assessment dated exactly \texttt{2026-08-21}: the supplied status is ``open''. This is supplied dated curation derived from the Oberwolfach Reports material, not an independently verified assessment. The supplied background expressly requests independent verification of time-sensitive status and cited claims.
\item The registry labels the statement \texttt{corrected\_verified} and reports checking its mathematical content against the official Oberwolfach report, with notation, context, and question form normalized. This is supplied verification metadata; the present draft does not independently verify that claim or the external references.
\end{itemize}
\end{document}
